
\paragraph{Summary of findings}
\begin{itemize}
  \item \textbf{RQ1:} the headline performance is reproduced: F1 0.98 on the paper's protocol (paper: 0.97--0.99).
  \item \textbf{RQ2:} the component claims are not: GAT $\approx$ GCN $\approx$ GIN, embedding dimension and context
        window barely matter, and normalisation/risk weights do not significantly help (paired tests, $p \geq 0.07$).
  \item \textbf{RQ3:} 77\% of test webshells leak via near-duplicates (F1 $-$2 points without them); a token
        TF-IDF + LR baseline is better than every GNN, and graph statistics alone reach F1 0.92.
  \item \textbf{RQ4:} on four unseen benign projects the FPR rises from 0.4\% to 1.7\%; recall on webshells from
        new collections is 0.54--0.64, and only 6\% at a 0.1\% alarm budget; 10 of 12 small legitimate scripts are flagged. The detector has learned
        ``not framework code'' rather than ``malicious''.
  \item \textbf{Improvement:} training on five small benign applications lowers the GAT FPR from 2.5\% to 0.9\%
        on all benign test files ($p=0.006$, recall unchanged), but 8 of the 12 scripts remain flagged.
\end{itemize}
